\subsection*{7. Pursuing Stacks}
The seventh section opens on p.~42 with the work begun in March of the
previous year on the ``fondements de l'algèbre (co)homologique non
commutative'', also called ``algèbre homotopique''. On p.~43 Grothendieck
describes the intended two-volume \emph{A la Poursuite des Champs}, its
successive expansions, and its place at the beginning of
\emph{Réflexions Mathématiques}, already announced in \S1. His assessment
first places this work somewhat at the margin of the newer interests
described in the preceding sections, as a return to questions from before
1970. He names Verdier, Giraud, and Illusie, acknowledges Quillen's model
categories, and on p.~44 explains the internal structures and
noncommutative coefficients that remain the focus of his own inquiry.
The correspondence with Breen in 1975 leads into the proposed role of
$\infty$-groupoids as models of homotopy types, with truncated types for
n-groupoids; the subsequent correspondence with Brown brings him back
to this problem. The account records two changes on p.~44: a broader
notion initially called an $\infty$-groupoid is renamed ``champ en
$\infty$-groupoïdes'', or ``champ'' over the punctual topos, and
\emph{Histoire de Modèles} becomes a large, unforeseen digression before
the intended treatment of stacks. The relation to \S6 is through these
homotopy types and higher groupoids; p.~45 also explicitly reconnects the
work with the moduli multiplicities and their stratifications from
\S\S2 and 5, through a higher-dimensional van Kampen question. The same
page returns to the 1975--76 reflections on a ``complexe de De Rham à
puissances divisées'' and to ``schématisation'' of homotopy types, while
qualifying the corresponding chapter as particularly hypothetical.
The final paragraph revises the image of merely paying a debt to the
1955--70 period: unforeseen developments, represented by the nested tales
of the \emph{Thousand and One Nights}, account for the present work's
unexpected scale. The archive catalogue names
\href{https://grothendieck.umontpellier.fr/archives-grothendieck/}{\emph{Pursuing stacks. The Modelizing Story} (cotes 134-1--134-8) and \emph{Complexe de De Rham à puissance divisée} (cote 136)};
among later literature one may name Denis-Charles Cisinski's
\href{https://smf.emath.fr/publications/les-prefaisceaux-comme-modeles-des-types-dhomotopie}{\emph{Les préfaisceaux comme modèles des types d'homotopie}}.
(\href{https://grothendieck.umontpellier.fr/web/viewer.html?file=../119.pdf\#page=45}{Copy 1, pp.~42--45, beginning at site 44};
\href{https://grothendieck.umontpellier.fr/web/viewer.html?file=../119.pdf\#page=47}{renaming and digression on p.~44, site 46};
\href{https://grothendieck.umontpellier.fr/web/viewer.html?file=../119.pdf\#page=48}{connections and closing reassessment on p.~45, site 47}.)
